\section{How is a leak detected?}
\label{app:labels}

A rule on the tool log finds the first call that writes the secret out, by exact string match, so the main label needs no judge.
Checks of how leaks go out, and whether they are deliberate, are in \cref{app:leakchecks}.

\paragraph{The leak label.}
A call leaks if the protected value appears, in full, in its reply or in the arguments of any of its tool calls.
This covers channel messages, the end-of-turn reply that the harness posts to the channel, and file writes, wherever they write.
The value appearing in private reasoning does not count.
The label is the one used for gpt-oss-20b's test confirmation (\cref{app:confirmation}), applied unchanged to Qwen3-32B.
On Qwen3-32B, \qwenSourceWritebackCalls{} leaking calls, in \qwenSourceWritebackEpisodes{} episodes, write the value back into the protected file itself, one in the development episodes and one in the test episodes.
Excluding the development one changes the per-call result very little (\cref{app:robustness}).

\paragraph{The label in the Qwen3-32B round-level test.}
The plan for the Qwen3-32B test episodes fixed a slightly different leak.
There, a leak is the complete value posted to the team channel, or written into another file in the five scenarios other than the main one.
Scored calls stop before the secret is first written out anywhere.
We report that test as it was planned.
On gpt-oss-20b, the round-level analyses in \cref{app:gptoss-round} use a channel post as the leak, again with scored calls stopping before the secret is first written out.
